\documentclass[12pt,a4paper]{article}

%\usepackage[backend=biber, style=numeric-comp, sorting=none]{biblatex}
%\addbibresource{refs.bib}

%\usepackage{graphicx, csquotes, tikz, pgfplots, wrapfig}
%\pgfplotsset{compat = newest}
%\graphicspath{{./pictures/}}

\usepackage[dutch]{babel}
\usepackage{csquotes}
\usepackage{amsmath, amssymb}

\usepackage{listings}

\lstdefinestyle{mystyle}{
    numbers=left,                    
    numbersep=5pt,                  
}
\lstset{style=mystyle}

\numberwithin{equation}{section}

\usepackage[hidelinks]{hyperref}

\binoppenalty=100000
\relpenalty=100000

\allowdisplaybreaks

\title{Verslag Taak 1}
\author{Robin Elskens en Lina Martens}
\date{} %TODO: Submission date

\begin{document}

\maketitle
%
%TODO: Better
%
\section*{Opdracht 1}
\begin{align*}
  \intertext{We stellen $n\geq2$ en $r\in\mathbb{R}$. We zullen de hint volgen en $A_n(r)$ partieel integreren.}
  A_n(r) &= \int_{-1}^{1}\left(1-x^2\right)^n\cos\left(rx\right)\,dx \\
        &= \left[\left(1-x^2\right)^n\frac{1}{r}\sin\left(rx\right)\right]_{-1}^1 - \int_{-1}^{1}\left(-2n\right)x\left(1-x^2\right)^{n-1}\frac{1}{r}\sin\left(rx\right)\,dx 
\intertext{We kunnen in deze uitdrukking de eerste term schrappen aangezien $(1-x^2)$ in zowel $-1$ als $1$ nul is. We zullen nu de resterende term opnieuw partieel integreren met}
  u &= 2nx\left(1-x^2\right)^{n-1}\text{, zodat} \quad du = 2n\left(1-x^2\right)^{n-1} - 4n(n-1)x^2\left(1-x^2\right)^{n-2}\, dx \text{ en} \\
  dv &= \frac{1}{r}\sin\left(rx\right)\,dx \text{, zodat} \quad v = \frac{-1}{r^2}\cos\left(rx\right).\\
\intertext{We kunnen $du$ herschrijven als}
  du &= 2n\left(1-x^2\right)^{n-1} + 4n(n-1)\left(1-x^2\right)^{n-1} - 4n(n-1)\left(1-x^2\right)^{n - 2}\\
     &= 2n(2n-1)(1-x^2)^{n-1} - 4n(n-1)(1-x^2)^{n-2}.
\intertext{Dit geeft ons}
A_n(r) &= \left[2nx\left(1-x^2\right)^{n-1}\frac{-1}{r}\cos\left(rx\right)\right]_{-1}^1 \\
       &+ \frac{1}{r^2}\int_{-1}^{1}\left[2n(2n-1)(1-x^2)^{n-1} - 4n(n-1)(1-x^2)^{n-2}\right]\cos\left(rx\right)\, dx.
\intertext{Net als hierboven is de eerste term gelijk aan nul. De lineariteit van de integraal en de formule voor de rij $A_n(r)$ geven ons}
A_n(r) &= \frac{1}{r^2}\left[2n(2n-1)A_{n-1}(r)-4n(n-1)A_{n-2}(r)\right],
\end{align*}
wat te bewijzen was.
%
%TODO: Overfull hbox
%
\section*{Opdracht 2}
We gebruiken de volgende SymPy code om $A_0$ t.e.m. $A_{10}$ te bepalen\footnotemark{}.
\begin{lstlisting}[language=Python]
import sympy as sp
r = sp.symbols('r', nonzero=True)
x = sp.symbols('x')
A = [None]*11
A[0] = sp.integrate((1-x**2)**0*sp.cos(r*x), (x, -1, 1))
A[1] = sp.integrate((1-x**2)**1*sp.cos(r*x), (x, -1, 1)).simplify()
for n in range(2, 11):
    A[n] = (1/r**2)*(2*n*(2*n-1)*A[n-1]-4*n*(n-1)*A[n-2])
    A[n] = A[n].simplify()
for f in A:
    display(sp.collect(f, (sp.cos(r), sp.sin(r)))) 
\end{lstlisting}
Dit zijn de uitdrukkingen die we dan bekomen:
\begin{align*}
  A_{0}(r) &=\frac{2 \sin{\left(r \right)}}{r} \\
  A_{1}(r) &=\frac{4 \left( \sin{\left(r \right)} - r \cos{\left(r \right)} \right)}{r^{3}} \\
  A_{2}(r) &=\frac{16 \left( \left(3 - r^{2}\right) \sin{\left(r \right)}- 3 r \cos{\left(r \right)} \right)}{r^{5}} \\
  A_{3}(r) &=\frac{96 \left(\left(15 - 6 r^{2}\right) \sin{\left(r \right)} + \left(r^{3} - 15 r\right) \cos{\left(r \right)}\right)}{r^{7}} \\
  A_{4}(r) &=\frac{768 \left(\left(r^{4} - 45 r^{2} + 105\right) \sin{\left(r \right) + \left(10 r^{3} - 105 r\right) \cos{\left(r \right)}}\right)}{r^{9}} \\
  A_{5}(r) &=\frac{7680 \left(\left(15 r^{4} - 420 r^{2} + 945\right) \sin{\left(r \right)} + \left(- r^{5} + 105 r^{3} - 945 r\right) \cos{\left(r \right)}\right)}{r^{11}} \\
  A_{6}(r) &=\frac{92160 \left( \left(- r^{6} + 210 r^{4} - 4725 r^{2} + 10395\right) \sin{\left(r \right)} + \left(- 21 r^{5} + 1260 r^{3} - 10395 r\right) \cos{\left(r \right)}\right)}{r^{13}} \\
  A_{7}(r) &=\frac{1290240}{r^{15}} \bigl(\left(- 28 r^{6} + 3150 r^{4} - 62370 r^{2} + 135135\right) \sin{\left(r \right)}\\ &+ \left(r^{7} - 378 r^{5} + 17325 r^{3} - 135135 r\right) \cos{\left(r \right)}\bigr) \\
  A_{8}(r) &=\frac{20643840}{r^{17}} \bigl(\left(r^{8} - 630 r^{6} + 51975 r^{4} - 945945 r^{2} + 2027025\right) \sin{\left(r \right)} \\ &+ \left(36 r^{7} - 6930 r^{5} + 270270 r^{3} - 2027025 r\right) \cos{\left(r \right)} \bigr) \\
  A_{9}(r) &=\frac{371589120}{r^{19}} \bigl(\left(45 r^{8} - 13860 r^{6} + 945945 r^{4} - 16216200 r^{2} + 34459425\right) \sin{\left(r \right)} \\ &+ \left(- r^{9} + 990 r^{7} - 135135 r^{5} + 4729725 r^{3} - 34459425 r\right) \cos{\left(r \right)}\bigr) \\
  A_{10}(r) &=\frac{7431782400}{r^{21}} \bigl( \left(- r^{10} + 1485 r^{8} - 315315 r^{6} + 18918900 r^{4} - 310134825 r^{2} + 654729075\right) \sin{\left(r \right)} \\ &+ \left(- 55 r^{9} + 25740 r^{7} - 2837835 r^{5} + 91891800 r^{3} - 654729075 r\right) \cos{\left(r \right)}\bigr)
\end{align*}
\footnotetext{Lijn 11 van de code is afkomstig van \url{https://stackoverflow.com/a/16239378}.}
\vspace{-10mm} % Weird, but works?
%
%TODO: Better general structure
%TODO: Not self contained
%
\setcounter{section}{3}
\section*{Opdracht 3}
We zullen de correctheid van $(1.4)$ aantonen met een inductie op $n$. Merk op dat we $A_0(r)$ en $A_1(r)$, als gevonden in de vorige opdracht, kunnen schrijven als
\begin{align*}
    A_0(r) &= \frac{0!}{r}\left(2\sin{\left(r \right)} + 0 \cos{\left(r \right)}\right)\\
    A_1(r) &= \frac{1!}{r^3}\left(4\sin{\left(r \right)} - 4r\cos{\left(r \right)}\right).
\end{align*}
%NOTE: Slordig...
Hiermee hebben we de basisstap aangetoond en hebben we ook al beginwaarden $P_0\left(r\right)= 2$ en $P_1\left(r\right)= 4$.
Zij nu $n\in\mathbb{N}$ met $n\geq2$ en veronderstel dat we $A_k$ kunnen schrijven als in (1.4), van zodra $k<n$. We tonen aan dat $A_n$ dan ook aan (1.4) voldoet.
Uit de recursieformule vinden we
\begin{align*}
    A_n(r) &= \frac{1}{r^2}\left[2n(2n-1)A_{n-1}(r)-4n(n-1)A_{n-2}(r)\right]\\
           &= \frac{1}{r^2}\left[2n(2n-1)\frac{\left(n-1\right)!}{r^{2n-1}}R_{n-1}(r)-4n(n-1)\frac{\left(n-2\right)!}{r^{2n-3}}R_{n-2}(r)\right]\\
           &= \frac{n!}{r^{2n+1}}\left[\left(4n-2\right)R_{n-1}(r)-4r^2R_{n-2}\right],
    \intertext{waarbij $R_i(r) := P_i(r)\sin{\left(r\right)} - Q_i(r)\cos{\left(r\right)}$. We concluderen dat}
    A_n(r) &= \frac{n!}{r^{2n+1}}\left[\left(\left(4n-2\right)P_{n-1}(r)-4r^2P_{n-2}(r)\right)\sin{\left(x\right)}-\left(\left(4n-2\right)Q_{n-1}(r)-4r^2Q_{n-2}(r)\right)\cos{\left(x\right)}\right].
\end{align*}
%TODO: overfull hbox
Hieruit volgt ook de recursierelatie 
\begin{equation}
P_n(r) = \left(4n-2\right)P_{n-1}(r)-4r^2P_{n-2}(r).
\label{eq:3.1}
\end{equation}

Uit deze formule vinden we inductief dat alle coëfficienten van $P_n$ gehele getallen zijn, aangezien de coëfficienten van $P_0$ en $P_1$ geheel zijn.
%TODO: Onvoldoende? 
% expressie AN opschonen
%
%
\subsubsection*{Graden van $P_i$ en $Q_i$}
Uit de bekomen recursierelaties voor $P_i(r)$ en $Q_i(r)$ staat er ons nog voor de graden van deze veeltermen te bepalen.
Hiervoor leggen we eerst de aandacht op de graad van $P_n(r)$, we beweren namelijk dat deze kan berekend worden door de formule:

\begin{equation}
\text{Graad$(P_n(r))=$}
\begin{cases}
i-1 & \text{als $i$ oneven is}\\
i & \text{als $i$ even is}
\end{cases}
\label{eq:3.2}
\end{equation}

Hierbij introduceren we ook een notatie die van pas zal komen om \ref{eq:3.2} mee aan te kunnen tonen.


\newtheorem{notation}{Notatie}[section]
\begin{notation}
Schrijf $a_n$ voor de leidende coëfficient van de veelterm $P_n$. Per keuze zal $a_n\neq0$ en definieer
$b_n=(4n-2)a_{n-1}-4a_{n-2}$.
\end{notation}

De formule kan nu met behulp van inductie afgeleid worden.

Beschouw de veeltermen:
$$P_0=2 \qquad \text{en} \qquad P_1=4$$
Er geldt voor $P_0$ en $P_1$ dat beiden graden $0$ zijn zodat formule \ref{eq:3.2} voor de beide veeltermen standhoud.
Veronderstel dus dat voor alle  $i \in\mathbb{N}$ met $0\leq i \leq n-1$,
\ref{eq:3.2} de juiste graad geeft voor $P_i$ en dat 

\begin{equation}
a_i=
\begin{cases}
-4a_{n-2} & \text{als $i$ oneven is}\\
(4n-2)a_{n-1}-4a_{n-2} & \text{als $i$ even is}
\end{cases}
\label{eq:3.3}
\end{equation}
En daar te boven wanneer $i$ oneven is, dat dan 

\begin{equation}
a_i=i(i+1)a_{i-1}.
\label{eq:3.4}
\end{equation}


%vind formule voor a_n

\vspace{2.5mm}
Om de inductie af te kunnen maken: hoeven we te bewijzen dat formule \ref{eq:3.2} de juiste graad van $P_n$ oplevert, dat \ref{eq:3.3} waar is, en in het geval dat $n$ oneven is, dat dan $a_n=n(n+1)a_{n-1}$. Aldus beschouwen we de gevallen waarin $n$ even is of $n$ oneven afzonderlijk.
\\
\begin{enumerate}
\item $n$ is even \\
Uit de formule \ref{eq:3.2} zal Graad$(P_{n-1}(r))=n-2$ en Graad$(P_{n-2}(r))=n-2$, dus volgens de recursiebetrekking \ref{eq:3.1} zal de graad van $P_n(r)$ bepaald worden door $-4r^2P_{n-2}(r)$, zodat Graad($P_n(r))=n$.
Daarmee zal ook $a_n=-4a_{n-2}$.
\item $n$ is oneven \\
Uit de formule \ref{eq:3.2} zal Graad$(P_{n-1}(r))=n-1$ en Graad$(P_{n-2}(r))=n-3$. We zullen hoeven voorzichtig om te gaan met de recursiebetrekking \ref{eq:3.1}, want de coëfficient van $P_n(r)$ die aan de macht $r^{n-1}$ zit is gelijk aan $b_n$ en er hoeft namelijk nog getoond te worden dat $b_n\neq0$ om te kunnen verzekeren dat weldegelijk Graad$(P_n(r))=n-1$. 

\begin{align*}
b_n
&=(4n-2)a_{n-1}-4a_{n-2}\\
&=(4n-2)a_{n-1}-4(n-2)(n-3)a_{n-3} \\
&=(4n-2)a_{n-1}+(n-2)(n-1)a_{n-1} \\
&=(n^2+n)a_{n-1} \\
&=n(n+1)a_{n-1} \\
\end{align*}


Hieruit is te ontleiden dat $b_n\neq0$ zodat $a_n=b_n$. De formules \ref{eq:3.3} en \ref{eq:3.4} blijven hiermee juist voor $n$ 
\end{enumerate}
Dit rond onze inductiehypothese af voor beide gevallen van $n$, \ref{eq:3.2} is daarmee bewezen.
%Geval Q nog uitwerken want bovenstaande redenering verlengt hiernaar niet onmiddelijk!
\\Voor de functies $Q_i(r)$ zijn 
$$Q_0=0 \qquad \text{en} \qquad Q_1=-4r$$
en kunnen wij gelijkaardig $c_n$ schrijven voor de leidende coëfficient van $Q_n(r)$ en $d_n$ voor de waarde $(4n-2)c_{n-1}-4c_{n-2}$

We beweren dat 
\vspace{1.5pt}
\begin{equation}
\text{Graad$(Q_n(r))=$}
\begin{cases}
i & \text{als $i$ oneven is}\\
i-1 & \text{als $i$ even is}
\end{cases}
\label{eq:3.5}
\end{equation}

Veronderstel uit inductie dat deze formule waar is voor $i$ met $0\leq i<n$.
In het geval dat $n$ even is zal de leidende coëfficient van $Q_n(r)$ overeen komen met $c=c_n$.
En voor $n$ oneven wordt de leidende coëfficient $-4c_{n-2}$. 
\vspace{2mm}

Dus kunnen we analoog door telkens in het bewijs van \ref{eq:3.2} het woord even met oneven te vervangen en vice versa, bewijzen dat \ref{eq:3.5} waar is.


\setcounter{section}{4}
\section*{Opdracht 4}
Zoals beschreven in de opgavenbundel is voor $r=\pi/2$ de integrand van $A_n(r) = \int_{-1}^{1}\left(1-x^2\right)^n\cos\left(rx\right)\,dx$ positief, dus is $0<A_n\left(\frac{\pi}{2}\right).$ 
Hieruit volgt dat het voldoende is om aan te tonen dat 
\begin{equation}
    \frac{a^{2n+1}}{n!}A_n\left(\frac{\pi}{2}\right)\in\mathbb{Z}.\label{eq:4.1}
\end{equation} 
Uit de vorige opgave volgt dat 
$$    \frac{a^{2n+1}}{n!}A_n\left(\frac{\pi}{2}\right) =\frac{a^{2n+1}}{n!}\frac{n!}{\left(\frac{a}{b}\right)^{2n+1}}\left[P_n\left(\frac{a}{b}\right)\sin{\left(\frac{a}{b}\right)}-Q_n\left(\frac{a}{b}\right)\cos{\left(\frac{a}{b}\right)}\right].$$
    Omdat $a/b = \pi/2$, kunnen we \eqref{eq:4.1} herleiden tot
   \begin{equation}
       b^{2n+1}P_n\left(\frac{a}{b}\right)\in\mathbb{Z}. \label{eq:4.2}
   \end{equation}
We tonen dit aan met inductie. In opdracht 3 zagen we dat $P_0(r) = 2$ en $P_1(r) =4$, dus de basisstap is voldaan. Zij nu $n\geq 2$ en veronderstel dat \eqref{eq:4.2} voldaan is voor elke $k<n$. We gebruiken de recursieformule voor $P_n$ zodat
\begin{align*}
     b^{2n+1}P_n\left(\frac{a}{b}\right)&=b^{2n+1}\left[\left(4n-2\right)P_{n-1}\left(\frac{a}{b}\right)-4\frac{a^2}{b^2}P_{n-2}\left(\frac{a}{b}\right)\right]\\
               &=\left(4n-2\right)b^2\left[b^{2n-1}P_{n-1}\left(\frac{a}{b}\right)\right] - 4a^2b^2\left[b^{2n-3}P_{n-2}\left(\frac{a}{b}\right)\right].
\end{align*}
Uit de inductiehypothese en de definitie van $a$ en $b$ volgt dat alle termen in deze gelijkheid geheel zijn, dus is het resultaat opnieuw geheel.

\section*{Opdracht 5}
Neem de reeks $\{b_n\}_{n\in\mathbb{N}}=\{\frac{c^{n}}{n!}\}_{n\in\mathbb{N}}$ waarbij $c\in\mathbb{R}$. Kies een $N$ waarvoor $N>c$. Voor $n>N$ zal
$$b_{n+1}/b_n=\frac{c}{n}<1$$ zodat $b_{n+1}<b_n$ en dus de reeks $\{b_n\}_{n>N}$ monotonisch daalt. Hieruit volgt er dat $\{b_n\}_{n\in\mathbb{N}}$ convergeert. 
Laat dus $b=\lim_{n\to\infty}b_n$, aangezien dat $b_{n+1}=\frac{c}{n}b_n$: kunnen we door de limiet van beiden leden te nemen krijgen dat $b=0$.

aldus, kies $c=a^2:$ waaruit
$$\lim_{n\to\infty}\frac{n!}{a^{2n+1}}=\lim_{n\to\infty}\frac{1}{ab_n}=\infty$$


\end{document}
